% SPDX-FileCopyrightText: 2026 Will Cook
% SPDX-License-Identifier: CC-BY-4.0
\documentclass[11pt]{article}
% SPDX-FileCopyrightText: 2026 Will Cook
% SPDX-License-Identifier: CC-BY-4.0
%
% Shared preamble for the Erdős Problem Notes: the short, standalone,
% problem-owned manuscripts covering the expansion library `ErdosProblems`.
%
% One note owns one Erdős problem.  Each note repeats whatever context its own
% reader needs, so that a reader who arrives at a single problem never has to
% assemble it from the gateway paper.  Deliberate repetition across notes is the
% design, not an oversight.
%
% Source links resolve against one pinned formal-source commit, declared once
% below.  `scripts/check_problem_note_sources.py` verifies that every authored
% (file, line, declaration) triple names that exact declaration in the pinned
% snapshot, so a link cannot silently rot as later work moves lines.

\usepackage{paper-house-style}
\usepackage{array}
\usepackage{booktabs}
\usepackage{enumitem}
\setlist{topsep=0.35em,itemsep=0.2em}
\usepackage[colorlinks=true,linkcolor=linkinternal,urlcolor=linkexternal,
  citecolor=linkinternal]{hyperref}
\hypersetup{bookmarksnumbered=true,bookmarksopen=true,bookmarksopenlevel=1,
  pdfauthor={Will Cook},pdflang={en-GB}}

% ---------------------------------------------------------------------------
% Pinned formal source.  \commit names the checkpoint every link resolves
% against; the checker reads that snapshot rather than the working tree, so the
% printed coordinates stay correct however later waves move declarations.
% ---------------------------------------------------------------------------
\newcommand{\commit}{e5fd26a6d1b1a346430205eab5385b4c9565d004}
\newcommand{\commitshort}{e5fd26a6d1b1}
\newcommand{\repobase}{https://github.com/wcook04/plectis-lean-erdos249-257/blob/\commit}
\newcommand{\sourceurl}{https://github.com/wcook04/plectis-lean-erdos249-257/tree/\commit}
\newcommand{\PX}{ErdosProblems}

% Declaration names stay inside link targets.  The printed label is either a
% spaced, lower-case reading of the declaration name (\lref) or a short authored
% phrase (\lword); neither prints a module path or a raw Lean identifier.
\ExplSyntaxOn
\cs_new_protected:Npn \noteleanlabel #1
  {
    \tl_set:Nx \l_tmpa_tl { \tl_to_str:n {#1} }
    \regex_replace_all:nnN { _+ } { \c{space} } \l_tmpa_tl
    \regex_replace_all:nnN
      { ([a-z0-9])([A-Z]) }
      { \1\c{space}\2 }
      \l_tmpa_tl
    \tl_set:Nx \l_tmpa_tl { \tl_lower_case:n { \tl_use:N \l_tmpa_tl } }
    \tl_use:N \l_tmpa_tl
  }
\ExplSyntaxOff
\newcommand{\leanlabel}[1]{\noteleanlabel{#1}}
\newcommand{\lref}[3]{\href{\repobase/\PX/#1\#L#2}{\leanlabel{#3}}}
\newcommand{\lword}[4]{\href{\repobase/\PX/#1\#L#2}{#4}}
\newcommand{\lloc}[2]{\href{\repobase/\PX/#1\#L#2}{Lean source}}

% Links into a sibling library, named repository-relative.  The expansion
% library is implicit in \lref/\lword, because that is where problem-owned
% work lands.  Several problems have their headline theorem in the older
% reviewed #249/#257 corpus, and a note that could not name that library
% would have to paraphrase a checked statement instead of linking it.
\newcommand{\mref}[3]{\href{\repobase/#1\#L#2}{\leanlabel{#3}}}
\newcommand{\mword}[4]{\href{\repobase/#1\#L#2}{#4}}
\newcommand{\mloc}[2]{\href{\repobase/#1\#L#2}{Lean source}}
\emergencystretch=2em

\newtheorem{theorem}{Theorem}[section]
\newtheorem{proposition}[theorem]{Proposition}
\newtheorem{lemma}[theorem]{Lemma}
\newtheorem{corollary}[theorem]{Corollary}
\theoremstyle{definition}
\newtheorem{definition}[theorem]{Definition}
\newtheorem{example}[theorem]{Example}
\newtheorem{problem}[theorem]{Problem}
\theoremstyle{remark}
\newtheorem*{remark}{Remark}

\newcommand{\N}{\mathbb{N}}
\newcommand{\Npos}{\mathbb{N}_{>0}}
\newcommand{\Z}{\mathbb{Z}}
\newcommand{\Q}{\mathbb{Q}}
\newcommand{\R}{\mathbb{R}}
\newcommand{\lcm}{\operatorname{lcm}}
\newcommand{\ph}{\varphi}

% The series line printed under every note's title block.
\newcommand{\seriesline}{Erd\H{o}s Problem Notes}

% Production provenance.  The wording is fixed by the corpus provenance
% contract and reproduced verbatim in every manuscript in this repository, so
% the papers cannot drift into different accounts of how they were made.
\newcommand{\noteprovenance}{\provenancenote{\emph{Authorship and AI use.} The
prose, formal proofs, and software in this version were drafted and revised by
large-language-model agents under Will Cook's direction.  Cook set the
objectives and acceptance criteria, reviewed and authorised the manuscript, and
is responsible for its contents.  The agents are production tools, not authors.
See \emph{Statements and declarations}.}}

% The status paragraph every note carries, in its own words, before any result.
% Kernel checking establishes the exact Lean proposition; it does not establish
% that the proposition is the interesting one, that it is new, or that the
% Erdős problem is closed.
\newcommand{\statusboundary}{%
  \emph{Status.}  The problem treated here is open, and this note does not
  close it.  Every statement below marked as checked is a proposition that the
  pinned Lean kernel accepts from the sources this note links to, with no
  \texttt{sorry}, no added axiom, and no unchecked evaluation.  That is a claim
  about the formal statement, not about its mathematical interest, its
  novelty, or the original problem.  The unresolved obligations are named
  exactly, in their own section, and none of the finite computations,
  reductions, or no-go results here removes one of them.}

\renewcommand{\commit}{32545d77c4b5cfd72ebd36c8dcd2418bf7cc4d31}
\renewcommand{\commitshort}{32545d77c4b5}
\hypersetup{
  pdftitle={A basis for the 2-kernel of Euler's totient},
  pdfsubject={Erdős Problem Notes: Problem 249},
  pdfkeywords={irrationality, Euler totient, k-regular sequences, 2-kernel, Lambert series, Möbius inversion, linear independence, Lean 4},
}

\newcommand{\Stot}{S}
\newcommand{\EB}{E}
\newcommand{\Aw}{A}
\newcommand{\Ker}{\mathcal K}
\newtheorem{equivform}[theorem]{Equivalent formulation}

% Status tag from the claim taxonomy, printed on every numbered result.
\newcommand{\stag}[1]{{\normalfont\footnotesize\textsc{[#1]}}}

\runninghead{Erd\H{o}s \#249: a basis for the dyadic totient sections}
\title{A Basis for the 2-Kernel of Euler's Totient}
\subtitle{Exact finite-level ranks and consequences for Erd\H{o}s
Problem~\#249}
\author{Will Cook}
\date{July 2026}

\begin{document}
\maketitle
\noteprovenance

\begin{abstract}
For every $e\ge1$, the dyadic sections of Euler's totient through level $e$
have rational rank $2^e+1$. The two zero-residue sections and the odd-residue
sections form a basis, and two families of scalar reductions generate every relation.
Divisibility makes an evaluation matrix diagonal modulo a prime, with nonzero
diagonal entries. The construction gives rank $k^e+1$ in every base $k\ge2$
and an integral normal form.
Separately, $\sum_{n\ge1}(\ph(n)\bmod m)2^{-n}$ is irrational for every
$m\ge3$, with a complete rationality classification at dyadic moduli.
For $S=\sum_{n\ge1}\ph(n)2^{-n}$, one nonintegral tail shift forces
full-depth certificates in every sufficiently late pair of multipliers.
Irrationality is equivalent to an explicit central-residue gap for every
prospective denominator; its arithmetic supply remains open.
\end{abstract}



\section{A basis and all its relations}\label{sec:results}
Set $\ph(0)=0$. For $k\ge2$, the $k$-kernel consists of the sequences
$n\mapsto\ph(k^jn+r)$ with $j\ge0$ and $0\le r<k^j$.
The following theorem gives a unique coordinate system at every finite level.
\begin{theorem}[exact rank of every totient $k$-kernel truncation]
\label{thm:kkernelrank}
Let $k\ge2$ and $e\ge1$ be integers, write $F^{(k)}_{j,r}(n)=\ph(k^jn+r)$, and put
\[
 V_{k,e}=\operatorname{span}_{\Q}
 \{\,F^{(k)}_{j,r}:0\le j\le e,\ 0\le r<k^j\,\}.
\]
Then $\dim_{\Q}V_{k,e}=k^e+1$, and
\[
 \mathcal B_{k,e}
 =\{F^{(k)}_{0,0},F^{(k)}_{1,0}\}
 \cup\{\,F^{(k)}_{j,r}:1\le j\le e,\ 1\le r<k^j,\ k\nmid r\,\}
\]
is a basis.  Every omitted section reduces to a basis element by an explicit
scalar: $F^{(k)}_{j,0}=k^{j-1}F^{(k)}_{1,0}$ for $j\ge1$, and if $r=k^tu$ with
$t=\max\{s:k^s\mid r\}\ge1$ and $k\nmid u$, then
\[
 F^{(k)}_{j,r}=C_k(t,u)\,F^{(k)}_{j-t,u},
 \qquad
 C_k(t,u)=k^t\prod_{\substack{p\mid k\\ p\nmid u}}\Bigl(1-\tfrac1p\Bigr).
\]
\end{theorem}
\label{res:rank}
The retained condition is $k\nmid r$. At $k=6$, for example, the residue
$r=2$ is retained. At base $2$ and level $3$, the fifteen sections have nine
coordinates; the omitted positive-residue sections satisfy
\[
 \ph(8n+2)=\ph(4n+1),\qquad
 \ph(8n+4)=2\ph(2n+1),\qquad
 \ph(8n+6)=\ph(4n+3).
\]

\subsection{The evaluation matrix}\label{sec:rank}
The two ingredients are scalar reduction and arithmetic independence.
The local formula for $\ph$ gives the displayed reductions. For a prime
$p\mid k$ that does not divide $u$, the argument $k^{j-t}n+u$ is a $p$-adic
unit, so multiplication by $k^t$ contributes its power of $p$ and one factor
$1-1/p$. When $p\mid u$, the argument already contains $p$, and only its
power contributes. This gives $C_k(t,u)$.

For independence, consider finitely many forms $L_i(n)=a_in+b_i$ with
positive slopes and $a_ib_j\ne a_jb_i$ for $i\ne j$. A finite shift makes
the intercepts positive. Write $L_i=g_iA_i$, where $A_i$ is primitive.
Choose an odd prime $\ell$ avoiding the slopes, all $g_i\ph(g_i)$, and the
nonzero cross determinants. In row $i$, choose distinct primes
$q_{ij}\equiv1\pmod\ell$ for $j\ne i$, outside the same finite exceptional
set, and impose
\[
 L_j(n)\equiv0\pmod{q_{ij}},\qquad A_i(n)\equiv2\pmod\ell.
\]
The Chinese remainder theorem combines these congruences. The resulting
progression for $A_i(n)$ is reduced: a prime $q_{ij}$ dividing $A_i(n)$
would also divide the nonzero cross determinant; $\ell$ is excluded by the
chosen residue; and primitivity excludes the primes of the slope.
Dirichlet's theorem supplies arbitrarily large primes $p=A_i(n)$ with
$p\nmid g_i$. Only the selected affine form is made prime in this row.
The diagonal value is
$\ph(g_i)(p-1)\not\equiv0\pmod\ell$. Each off-diagonal value is divisible
by $\ell$, since $q_{ij}-1$ divides the corresponding totient.
One evaluation per row gives
\[
 \bigl(\ph(L_j(n_i))\bigr)_{i,j}
 \equiv\operatorname{diag}(u_1,\ldots,u_s)\pmod\ell,
 \qquad u_i\ne0.
\]
The determinant is nonzero over $\Q$.

Apply this construction after restricting a relation among the retained
sections to $n=km+1$. The positive-residue sections give pairwise
nonproportional forms
\[
 k^{j+1}m+(k^j+r),\qquad k\nmid r,
\]
and the remaining form is $km+1$. Equality between two affine fractions
would make the residue at the larger level divisible by $k$. The matrix
therefore eliminates every positive-residue coefficient. If $A,B$ are the
two zero-residue coefficients, the restriction gives
$A+B\ph(k)=0$. Evaluation at $n=k$ gives $A+Bk=0$, since
$\ph(k^2)=k\ph(k)$. As $\ph(k)<k$, both coefficients vanish.
The reductions give spanning and the count is
\[
 2+\sum_{j=1}^e(k^j-k^{j-1})=k^e+1.
\]
This proves Theorem~\ref{thm:kkernelrank}. The scalar reduction at level
$j-t$, with the Euler product written out, is
\lloc{Erdos249/PaperCompleteR7/KernelIntegral.lean}{226}.

\begin{corollary}[Dyadic basis]\label{res:basis}
The family
\[
 \{\ph_{0,0},\ph_{1,0}\}
 \cup\{\ph_{j,r}:j\ge1,\ 0<r<2^j,\ r\text{ odd}\},
 \qquad \ph_{j,r}(n)=\ph(2^jn+r),
\]
is a basis for the rational span of the full dyadic kernel. Every rational
relation is generated by the scalar reductions. The complete level-zero
truncation has dimension one.
\end{corollary}
\begin{proof}
Every finite subfamily lies in a sufficiently high truncation. The theorem
gives its independence; the reductions give equality of the full spans.
The level-zero clause is
\lloc{Erdos249/PaperCompleteR7/ArithmeticAssemblies.lean}{112}.
\end{proof}

\begin{corollary}[Integral normal form]\label{cor:integral-normal-form}
The canonical family is also a $\Z$-basis of the module generated by the
sections through level $e$. Its integral relation module has a basis with
one elementary reduction per omitted channel and rank
$\sum_{j=1}^{e-1}k^j$.
\end{corollary}
\begin{proof}
The denominator $\prod_{p\mid k,\,p\nmid u}p$ in $C_k(t,u)$ divides $k$,
so all reduction scalars are integers. In the free module on the channels,
subtracting one elementary reduction for each omitted coordinate leaves a
relation among the independent retained coordinates. The omitted-coordinate
pivots are all one, which also gives independence of the elementary relations.
The coordinate module carries
\lword{Erdos249/PaperCompleteR7/KernelIntegral.lean}{138}{integralTotientKernelBasis}{an integral basis}
in Lean. The basis of the relation module stays outside the Lean development.
\end{proof}
The module here is the one generated by the channels. It is not saturated
in the group of all integer-valued sequences: at base $2$ and $e\ge2$,
$\ph(4n+3)/2$ is integer-valued but has canonical coefficient $1/2$.

\subsection{A reusable consequence}
\begin{corollary}[Freezing periodic coefficients]\label{cor:periodic-freezing}
If $L_i$ are pairwise nonproportional positive-slope affine forms and $w_i$
are rational-valued periodic sequences, then
\[
 \sum_i w_i(n)\ph(L_i(n))=0\quad\text{for all sufficiently large }n
 \quad\Longrightarrow\quad w_i(n)=0\quad\text{for every }i,n.
\]
\end{corollary}
\begin{proof}
Choose a common period $Q$. On each progression $n=a+Qt$, the coefficients
are constant and the cross determinants become $Q(a_ib_j-a_jb_i)\ne0$.
Shift beyond the threshold and apply the affine independence just proved.
The statement is \lloc{Erdos249/PaperCompleteR7/PeriodicAndPulse.lean}{64},
with a separate period for each coefficient.
\end{proof}
The same argument preserves the basis after deletion of a finite prefix.
For the two zero channels, large $n\equiv1\pmod k$ and large powers of $k$
replace the two evaluations used above.
\section{Bounded residues and rationality}\label{sec:family}
The residue theorem uses the same separation principle: force neighbouring
coefficients to vanish while retaining one central value.

\begin{theorem}\label{res:residueseries}
For every $m\ge3$,
\[
 A_m=\sum_{n\ge1}\frac{\ph(n)\bmod m}{2^n}\notin\Q.
\]
For $k\ge1$ and $f:\Z/2^k\Z\to\Q$, the series
$\sum_{n\ge1}f(\ph(n)\bmod2^k)2^{-n}$ is rational exactly when $f$ is
constant on the even residue classes. If that constant is $c$, its value is
$3f(1)/4+c/4$.
\end{theorem}

\begin{lemma}[A bounded isolated pulse]\label{lem:bounded-pulse}
Let $a_n\in\Z$ satisfy $|a_n|\le C$. Suppose that for arbitrarily large $L$
there is $N>L$ such that $a_N\ne0$ and
$a_{N+t}=0$ for $0<|t|\le L$. Then $\sum_{n\ge1}a_n2^{-n}$ is irrational.
\end{lemma}
\begin{proof}
If the sum has denominator $q$, every scaled tail
$T_j=\sum_{t\ge1}a_{j+t}2^{-t}$ belongs to $q^{-1}\Z$.
Choose $L$ so large that $C2^{-L}<1/q$.
The right zero block gives $|T_N|\le C2^{-L}<1/q$, hence $T_N=0$.
The left zero block now gives
$T_{N-L-1}=a_N2^{-L-1}$, a nonzero element of $q^{-1}\Z$ with absolute
value less than $1/q$, a contradiction. The statement is
\lloc{Erdos249/PaperCompleteR7/PeriodicAndPulse.lean}{106}.
\end{proof}

\begin{proof}[Proof of Theorem~\ref{res:residueseries}]
Fix a unit $s\pmod m$ and a block length $L$. For each nonzero
$t\in[-L,L]$, choose a distinct prime $q_t\equiv1\pmod m$ larger than $L$
and outside the prime divisors of $m$. Prescribe
\[
 p\equiv s\pmod m,\qquad p\equiv-t\pmod{q_t}.
\]
This is a reduced progression, because each $q_t>|t|$. Dirichlet supplies
arbitrarily large primes $p$ in it. Then $m\mid\ph(p+t)$ for every
$0<|t|\le L$, whereas $\ph(p)\equiv s-1\pmod m$.

Subtract $f(0)$ from the observable and clear its denominators. Whenever
$f(s-1)\ne f(0)$, the preceding construction gives arbitrarily long bounded
isolated pulses, so Lemma~\ref{lem:bounded-pulse} applies. For the
least-residue map choose $s=-1$: its central residue is $m-2\ne0$ when
$m\ge3$. For $m=2^k$, every even residue $r$ has $r+1$ a unit, so a
nonconstant restriction to the even residues supplies such a pulse.
Conversely, $\ph(n)$ is even for every $n\ge3$. Constancy on the even
classes therefore makes all terms after $n=2$ constant, giving the stated
rational value. This converse requires no recurrent-support theorem.
The three clauses together are
\lloc{Erdos249/PaperCompleteR7/RationalObservableClassification.lean}{277}.
\end{proof}
The bounded-pulse proof fixes its coefficient bound before choosing the block.
For $S$, the termwise condition $2^L\mid\ph(N+L)$, with $N+L>1$, forces
$2^L\le\ph(N+L)\le N+L-1$. Hence the tail envelope
$(N+L+2)2^{-L}>1$ cannot exclude any positive denominator.
The following criterion instead measures cancellation in a weighted prefix.

\section{The actual series and the remaining residue gap}\label{sec:carry-rank}
Write
\[
 S=\sum_{n\ge1}\ph(n)2^{-n},\qquad
 R_N=\sum_{j\ge1}\ph(N+j)2^{-j},\qquad
 \Delta_h(N)=R_{N+h}-R_N.
\]
The prefix identity $2^NS=\Phi_N+R_N$, with
$\Phi_N=\sum_{n=1}^N\ph(n)2^{N-n}\in\Z$, gives
\begin{equation}\label{eq:tail-phase}
 \Delta_h(N)\equiv2^N(2^h-1)S\pmod\Z.
\end{equation}
In particular a nonintegral tail difference is a precise arithmetic
statement about $S$.

\subsection{One seed removes the full-depth restriction}
For $h,N,L\in\N$, let
\[
 D_{h,N,L}=\sum_{j=0}^{L-1}
   \bigl(\ph(N+h+1+j)-\ph(N+1+j)\bigr)2^{L-1-j},
 \qquad B=N+h+L+2.
\]
Write $K(h,N,L)$ for
\[
 B<D_{h,N,L}\bmod2^L<2^L-B.
\]
The exact truncation and $\ph(n)\le n$ give
\begin{equation}\label{eq:tail-error}
 |2^L\Delta_h(N)-D_{h,N,L}|\le B.
\end{equation}
Thus some $L$ satisfies $K(h,N,L)$ exactly when $\Delta_h(N)\notin\Z$:
necessity follows from the strict inequalities; sufficiency follows by
choosing $L$ with $2^L\|\Delta_h(N)\|_{\R/\Z}>2B$.

\begin{theorem}[Full-depth amplification]\label{res:fulldepth}
Fix $d\ge1$ and $N\ge0$. If $\Delta_d(N)\notin\Z$, then every sufficiently
late pair $\{t,t+1\}$ contains an $m$ such that $K(md,N,md)$ holds.
Consequently
\[
 \exists t\ge1:\ K(td,N,td)
 \quad\Longleftrightarrow\quad\Delta_d(N)\notin\Z.
\]
Requiring this for every $d\ge1,N\ge0$ is equivalent to $S\notin\Q$.
\end{theorem}
\begin{proof}
Set $F_t=\Delta_{td}(N)$. Equation~\eqref{eq:tail-phase} yields
$F_{t+1}\equiv2^dF_t+F_1\pmod\Z$. If $K(td,N,td)$ fails,
\eqref{eq:tail-error} gives
\[
 \|F_t\|_{\R/\Z}\le\epsilon_t,
 \qquad \epsilon_t=2(N+2td+2)2^{-td}.
\]
Two adjacent failures would imply
$\|F_1\|\le2^d\epsilon_t+\epsilon_{t+1}$, impossible for all sufficiently
large $t$ because the left side is positive and the right side tends to zero.
Conversely, if $F_1$ is integral, then every $F_t$ is integral by the same
recurrence, so no certificate exists. Finally, \eqref{eq:tail-phase} is
nonintegral for every $d,N$ when $S$ is irrational. If
$S=a/(2^cv)$ with $v$ odd, take $N=c$ and $d=\ph(v)$ to make it integral.
\end{proof}
For instance $D_{1,12,16}=-143140$, with residue $53468$ modulo $2^{16}$
and $B=31$, supplies a seed. The amplifier produces infinitely many
full-depth certificates from this finite fact; it does not provide seeds
for all $d,N$.

\subsection{A single canonical endpoint}\label{sec:frontier}
For $c\ge0$, odd $v\ge1$ and a positive multiple $H$ of $\ph(v)$, put
\[
 M=\frac{2^H-1}{v},\qquad
 B_{H,c}=\sum_{j=0}^{H-1}\ph(c+1+j)2^{H-1-j},\qquad K=c+H+1.
\]
Euler's theorem makes $M$ integral.

\begin{equivform}[Canonical Mersenne gap]\label{res:canonicalmersenne}
The series $S$ is irrational if and only if, for every $c\ge0$ and odd
$v\ge1$, some positive multiple $H$ of $\ph(v)$ satisfies
\begin{equation}\label{eq:canonical-gap}
 K<(-B_{H,c})\bmod M<M-K.
\end{equation}
\end{equivform}
\begin{proof}
The identity
\[
 B_{H,c}=2^HR_c-R_{c+H}=M(vR_c)-\Delta_H(c)
\]
is exact. For $H\ge1$, $|\Delta_H(c)|<c+H+1=K$. Indeed $R_N\le N+1$
for $N\ge1$ and $0<R_0\le3/2$; the strict inequalities follow by subtracting
a positive tail. If $S=a/(2^cv)$, then $vR_c\in\Z$ and
$\Delta_H(c)\in\Z$. The residue of $-B_{H,c}$ is therefore within distance
less than $K$ of $0\pmod M$, so \eqref{eq:canonical-gap} fails.

If $S$ is irrational, fix $c,v$. Then $vR_c$ is nonintegral. As $H$ tends
to infinity through positive multiples of $\ph(v)$,
$\Delta_H(c)/M\to0$ and $K/M\to0$. Hence the fractional part of
$-B_{H,c}/M=-vR_c+\Delta_H(c)/M$ stays away from both endpoints of $[0,1]$,
which gives \eqref{eq:canonical-gap} for all sufficiently large such $H$.
\end{proof}
This proof identifies the equivalent missing quantifier. It does not supply
$H$ arithmetically without assuming the desired irrationality. The
equivalence, with the modulus written as the natural quotient $(2^H-1)/v$
and $v$ odd, is
\lloc{Erdos249/PaperCompleteR7/ArithmeticAssemblies.lean}{156}.
The finite residues also satisfy
\[
 \rho_{H,N+1,M}=
 \bigl(2\rho_{H,N,M}-\ph(N+H+1)+\ph(N+1)\bigr)\bmod M,
 \quad \rho_{H,N,M}=(-B_{H,N})\bmod M,
\]
which is an exact recurrence for searches directed at \eqref{eq:canonical-gap}.

\section{Information that does not force the gap}\label{sec:nogo}
The long record supplies a rational control with
$c(n)=\ph(n)$ for odd $n$, $|c(n)-\ph(n)|\le2$ for even $n$, $0\le c(n)\le n$,
and
\[
 \sum_{n\ge1}c(n)2^{-n}=\frac54.
\]
It has an integral tempered carry with dyadic rank at least $2^e-1$ at every
level. Thus the same lower bound for a hypothetical totient carry is
compatible with rationality. For one fixed prime $p$, the exact laws
$g(pn)=p g(n)$ when $p\mid n$ and $g(pn)=(p-1)g(n)$ otherwise, together
with $g-\ph=o(n)$, force $g=\ph$: a nonzero defect keeps a nonzero ratio
to the argument along $p^j n$. The control fails these laws; its full
coefficient rank is asserted only at measured finite depths.

A different obstruction concerns rational approximation. Put
\[
 \Theta_2=\sum_{d\ge1}\frac{\mu(d)}{(2^d-1)^2}=S-\frac12,
 \qquad
 Q(e,Y)=\frac{\bigl(\sum_{d=1}^{Y}\mu(d)/(2^d-1)^{e+2}\bigr)^2}
 {\sum_{d=1}^{Y}\mu(d)/(2^d-1)^{2e+2}}.
\]
The divisor-convolution identity gives the displayed value of $\Theta_2$,
as recorded by Fan~\cite{fan2026totient}. The following exact extremum is
stated in the long record as Theorem~\texttt{thm:rankonefloor}; its proof is
\lword{Erdos249/RankOneSharpFloor.lean}{534}{rankOneSubrankQuotient_eq_one_five_iff}{the sharp-floor argument}
in the linked source.

\begin{theorem}[Positive rank-one floor]\label{res:rankonefloor}
For $e\ge1$ and $Y\ge4$, the denominator of $Q(e,Y)$ is positive, and the
unique minimiser is $(e,Y)=(1,5)$. Every admissible quotient and every
nonempty finite positive weighted average of such quotients satisfies
\[
 Q-\Theta_2>\frac{21}{320}.
\]
The uniform bound $Q-\Theta_2>1/15$ is false already at $(1,5)$.
\end{theorem}
This theorem treats averages after the quotients are formed. Signed
coefficients, coupling before the quotient, and higher-rank kernels require
separate arguments. For any proposed rational approximants $p_j/q_j$ in
lowest terms, a sufficient arithmetic target is
\[
 0<q_j\left|\Theta_2-p_j/q_j\right|\longrightarrow0.
\]
Ordinary convergence and Hankel nonvanishing do not contain this denominator
estimate. The canonical gap \eqref{eq:canonical-gap} remains the equivalent
endpoint; the approximation estimate is one stronger sufficient route.

\section{Antecedents and proof record}\label{sec:open}
Allouche and Shallit introduced $k$-regular sequences~\cite{allouche-shallit}.
Coons proved that $\ph$ is not $k$-regular~\cite{coons}; Martin's stronger
affine ordering theorem implies the independence in
Section~\ref{sec:rank}~\cite{martin-phi-inequalities}. The basis theorem
gives the exact finite-level normal form. Wong~\cite{wong2015} proved
least-residue irrationality when the denominator base equals the modulus;
Section~\ref{sec:family} fixes base two and classifies dyadic observables.
The original question appears in Erd\H{o}s and Graham~\cite[p.~61]{erdosgraham1980}.

The exact all-base basis and rank, residue-series irrationality, full-depth
amplification, canonical equivalence and positive rank-one floor have
proofs in the supplied Lean source snapshot:
\mword{Erdos249257/AllBaseTotientKernel.lean}{1248}{allBaseTotientKernelBasis}{basis},
\lword{Erdos249/ResidueClassTotientSeries.lean}{595}{residue_series_irrational}{residues},
\lword{Erdos249/FullDepthRayAmplifier.lean}{364}{eventually_twoSyndetic_fullDepthKillMultipliers_of_seed}{amplification},
\lword{Erdos249/CyclotomicAnchoredKill.lean}{3256}{fullMersenneCanonicalBasepointResidueGapSupply_iff_irrational}{canonical gap}, and
\lword{Erdos249/RankOneSharpFloor.lean}{624}{rankOneSubrankQuotient_sub_theta_two_gt_twentyOne_div_threeTwenty}{rank-one floor}.
The historical source links retain \texttt{\commitshort}; the all-base
module requires a matching public source supplement.
The integral normal form has its ordinary proof above. The periodic transport
lemma is checked in \texttt{PeriodicTotientIndependence.lean}; the eventual
version follows by the displayed threshold shift. The rational control and
the termwise-window and prime-dilation boundaries are proved in
\texttt{ParityPerturbedRationalControl.lean} and
\texttt{PrefixValuationAndControlRigidity.lean}, respectively.
Finite evaluation ranks and the all-level carry lower bound remain distinct.
The long record retains the separate continued-fraction bounds
$q\ge2^{39989}$ and $q>10^{12038}$ under rationality, alongside the
Lean-checked Farey exclusion. Neither implies the universal gap.
\begin{thebibliography}{9}\small
\bibitem{allouche-shallit} J.-P.~Allouche and J.~Shallit,
\emph{The ring of $k$-regular sequences}, Theoret.\ Comput.\ Sci.\
\textbf{98} (1992), 163--197. doi:10.1016/0304-3975(92)90001-V.
\bibitem{coons} M.~Coons, \emph{(Non)Automaticity of number theoretic
functions}, J.\ Th\'eor.\ Nombres Bordeaux \textbf{22} (2010), 339--352,
Theorem~3.2. doi:10.5802/jtnb.718.
\bibitem{martin-phi-inequalities} G.~Martin, \emph{Simultaneous inequalities
among values of the Euler phi-function}, 2006, arXiv:math/0603053,
Theorem~1 and Corollary~4.
\bibitem{wong2015} E.~Wong, answer to
\href{https://math.stackexchange.com/questions/1210886}{question 1210886},
Mathematics Stack Exchange, 29 March 2015.
\bibitem{erdosgraham1980} P.~Erd\H{o}s and R.~L.~Graham,
\emph{Old and New Problems and Results in Combinatorial Number Theory},
1980, p.~61.
\bibitem{fan2026totient} S.~Fan, comment on Erd\H{o}s Problem \#249,
16 May 2026, 19:01, \href{https://www.erdosproblems.com/forum/thread/249}
{Erd\H{o}s Problems discussion thread}.
\end{thebibliography}

\end{document}
